\documentclass[tikz,border=12pt]{standalone}
\usepackage[T1]{fontenc}
\usepackage{mathpazo}
\usepackage{amsmath,amssymb}
\usetikzlibrary{arrows.meta, positioning, calc, decorations.pathreplacing}
\definecolor{stblue}{HTML}{7B97AD}
\definecolor{rwgold}{HTML}{B8995E}
\definecolor{sage}{HTML}{7A9A7E}
\definecolor{dkbrown}{HTML}{3D3530}
\definecolor{wmgray}{HTML}{7B7067}
\definecolor{cream}{HTML}{F9F6F0}
\definecolor{border}{HTML}{D4C9B8}
\definecolor{terracotta}{HTML}{C47A6A}
\definecolor{lavender}{HTML}{907EA0}

\begin{document}
\begin{tikzpicture}[
  background rectangle/.style={fill=cream},
  show background rectangle,
  every node/.style={text=dkbrown, font=\small},
  vecstyle/.style={-{Stealth[length=5pt]}, line width=1.1pt},
]

% ============================================
% LEFT PANEL: Projected Gradient Descent
% ============================================
\begin{scope}[shift={(0,0)}]

  % Title
  \node[font=\normalsize\bfseries, text=stblue!80!black] at (2.5,4.8) {Projected Gradient Descent};

  % Feasible polytope C — a compact convex shape
  \coordinate (v1) at (0.8,0.8);
  \coordinate (v2) at (3.8,0.5);
  \coordinate (v3) at (4.5,2.2);
  \coordinate (v4) at (3.5,3.8);
  \coordinate (v5) at (1.2,3.2);

  \fill[stblue!12] (v1) -- (v2) -- (v3) -- (v4) -- (v5) -- cycle;
  \draw[stblue, line width=0.9pt] (v1) -- (v2) -- (v3) -- (v4) -- (v5) -- cycle;
  \node[font=\footnotesize, text=stblue!60!black] at (2.5,1.6) {$\mathcal{C}$};

  % Current point x^k inside feasible set
  \coordinate (xk) at (2.8,2.5);
  \fill[dkbrown] (xk) circle (3pt);
  \node[font=\footnotesize, above left=-1pt] at (xk) {$x^k$};

  % Gradient step goes OUTSIDE the feasible set (to the right and down)
  \coordinate (grad_target) at (4.8,1.0);
  \draw[vecstyle, terracotta!80!black] (xk) -- (grad_target);
  \fill[terracotta!60] (grad_target) circle (3pt);
  \node[font=\footnotesize, text=terracotta!70!black, right] at (grad_target) {$x^k - \alpha\nabla f(x^k)$};

  % Projection: closest point on boundary back to C
  % Project onto edge v2--v3, roughly at (4.15,1.35)
  \coordinate (proj_point) at (4.1,1.3);
  \draw[vecstyle, sage!80!black, densely dashed] (grad_target) -- (proj_point);
  \fill[sage] (proj_point) circle (3.5pt);
  \node[font=\footnotesize, text=sage!70!black, above left=-1pt] at ($(proj_point)+(0.1,0.25)$) {$x^{k+1}$};

  % Labels along arrows
  \node[font=\scriptsize, text=terracotta!70!black, rotate=-37] at ($(xk)!0.5!(grad_target)+(0.0,0.25)$) {gradient step};
  \node[font=\scriptsize, text=sage!70!black, rotate=25] at ($(grad_target)!0.5!(proj_point)+(0.05,0.2)$) {project};

  % Right-angle mark at projection
  \coordinate (mid_edge) at ($(v2)!0.57!(v3)$);

\end{scope}

% ============================================
% RIGHT PANEL: Frank-Wolfe
% ============================================
\begin{scope}[shift={(7.5,0)}]

  % Title
  \node[font=\normalsize\bfseries, text=lavender!80!black] at (2.5,4.8) {Frank-Wolfe};

  % Same feasible polytope C
  \coordinate (v1) at (0.8,0.8);
  \coordinate (v2) at (3.8,0.5);
  \coordinate (v3) at (4.5,2.2);
  \coordinate (v4) at (3.5,3.8);
  \coordinate (v5) at (1.2,3.2);

  \fill[lavender!10] (v1) -- (v2) -- (v3) -- (v4) -- (v5) -- cycle;
  \draw[lavender, line width=0.9pt] (v1) -- (v2) -- (v3) -- (v4) -- (v5) -- cycle;
  \node[font=\footnotesize, text=lavender!60!black] at (2.5,1.6) {$\mathcal{C}$};

  % Mark all vertices
  \foreach \p in {(v1),(v2),(v3),(v4),(v5)}{
    \fill[wmgray] \p circle (2pt);
  }

  % Current point x^k
  \coordinate (xk) at (2.8,2.5);
  \fill[dkbrown] (xk) circle (3pt);
  \node[font=\footnotesize, above right=-1pt] at (xk) {$x^k$};

  % Gradient arrow (showing direction, short)
  \draw[vecstyle, wmgray!70] (xk) -- ++(0.8,0.6);
  \node[font=\tiny, text=wmgray] at ($(xk)+(1.1,0.7)$) {$\nabla f(x^k)$};

  % Linear minimizer y^k: the vertex that minimizes <grad f, y>
  % Gradient points upper-right, so minimizer is lower-left vertex
  \fill[rwgold] (v1) circle (3.5pt);
  \node[font=\footnotesize, text=rwgold!80!black, below] at ($(v1)+(0,-0.1)$)
    {$y^k$};

  % Dashed line from x^k to y^k (Frank-Wolfe direction)
  \draw[terracotta!80!black, line width=1pt, densely dashed] (xk) -- (v1);

  % Convex combination point x^{k+1} along the direction
  \coordinate (xnew) at ($(xk)!0.35!(v1)$);
  \fill[sage] (xnew) circle (3.5pt);
  \node[font=\footnotesize, text=sage!70!black, above left=-1pt] at (xnew) {$x^{k+1}$};

  % Label
  \node[font=\scriptsize, text=terracotta!60!black, rotate=-40] at ($(xk)!0.65!(v1)+(-0.1,0.2)$)
    {FW direction};

\end{scope}

% ============================================
% Bottom comparison note
% ============================================
\node[draw=border, fill=cream, rounded corners=4pt, inner sep=6pt,
  font=\footnotesize, text=dkbrown, align=center]
  at (5,-1.2) {%
  PGD: gradient step $\to$ project onto $\mathcal{C}$ (needs \textbf{projection oracle})
  \hspace{2em}
  FW: linearize $\to$ minimize over $\mathcal{C}$ $\to$ convex combination (needs \textbf{linear oracle})
  };

\end{tikzpicture}
\end{document}
